\section{Definitions and Problem Statement}

\subsection{Definitions}
In this section, we introduce the basic concepts that will be used throughout
this paper. Let $\Sigma$ denote an alphabet, which is a finite set of
characters. A string $s$ is a sequence of characters derived from the alphabet
$\Sigma$. The length of a string $s$ is denoted as $|s|$.

\begin{definition}[Prefix and Suffix \cite{crochemore1994text}]
A prefix of a string $s = a_1 a_2 \dots a_n$ of length $j$ is a string denoted
by $s[:j+1] = a_1 a_2 \dots a_j$. Similarly, a suffix is a string denoted by
$s[i:] = a_i \dots a_n$.
\end{definition}

\begin{definition}[Substring \cite{crochemore1994text}]
A string $x$ is a substring of a string $s$ if there exist strings $u$ and $v$
(which may be empty) such that $s = uxv$. A set of strings $S = \{s_1, s_2,
\dots, s_m\}$ is considered substring-free if no string $s_i \in S$ is a
substring of any other string $s_j \in S$.
\end{definition}

\begin{definition}[Overlap, Prefix, and Distance \cite{blum1994linear}]
For two strings $s$ and $t$, let $y$ be the longest string such that $s = xy$
and $t = yz$ for some non-empty strings $x$ and $z$. We define the following
concepts:
\begin{itemize}
    \item The \textbf{overlap} between $s$ and $t$, denoted as $ov(s, t)$, is
    the length of the string $y$, i.e., $ov(s,t) = |y|$. This notion can be
    naturally extended to cases where $t$ is a semi-infinite string, provided
    that $s$ remains finite.
    \item The \textbf{prefix} of $s$ with respect to $t$, denoted as
    $\textit{pref}(s, t)$, is defined as the string $x$. Therefore, $s$ can be
    uniquely decomposed as $s = \textit{pref}(s,t) y$.
    \item The \textbf{distance} from $s$ to $t$, denoted as $d(s, t)$, is
    defined as the length of this prefix, meaning $d(s, t) = |\textit{pref}(s,
    t)| = |x|$.
\end{itemize}
\end{definition}

\begin{definition}[Common Superstring \cite{blum1994linear}]
A common superstring for a set of strings $S = \{s_1, s_2, \dots, s_m\}$ is a
string $s$ such that every $s_i \in S$ is a substring of $s$.
\end{definition}

\begin{definition}[Superstring of a Sequence \cite{blum1994linear}]
For a given ordered list of strings $s_{i_1}, \dots, s_{i_r}$, the shortest
common superstring representing them, denoted as $\langle s_{i_1}, \dots,
s_{i_r} \rangle$, is defined as the concatenation of the prefixes and the final
string:
\begin{equation}
\langle s_{i_1}, \dots, s_{i_r} \rangle = \textit{pref}(s_{i_1}, s_{i_2}) \textit{pref}(s_{i_2}, s_{i_3}) \dots \textit{pref}(s_{i_{r-1}}, s_{i_r}) s_{i_r}
\end{equation}
\end{definition}

Any optimal common superstring for a finite set of strings $S = \{s_1, \dots,
s_m\}$ always corresponds to a sequential construction of the form $\langle
s_{i_1}, \dots, s_{i_m} \rangle$ for some optimal permutation of indices $i_1,
\dots, i_m$ from the set $\{1, \dots, m\}$. The length of such a shortest
superstring is denoted as $\textit{opt}(S)$.

We define $\textit{maxov}(S)$ as the maximum total overlap between adjacent
elements in the superstring. Formally, it is the maximum possible sum of
overlaps $\sum_{j=1}^{m-1} ov(s_{i_j}, s_{i_{j+1}})$ achieved over all
permutations of the input sequence. The size of the optimal solution, denoted as
$\textit{opt}(S)$, is strictly related to this value by the equation:
\begin{equation}
\textit{opt}(S) = \sum_{s_i \in S} |s_i| - \textit{maxov}(S)
\end{equation}
From the above identity, it follows that the problem of minimizing the
superstring length $\textit{opt}(S)$ is equivalent to the problem of maximizing
the sum of overlaps $\textit{maxov}(S)$.

\subsection{Problem Statement and Applications}

\begin{problem}[Shortest Common Superstring]
Given a set of strings $S = \{s_1, s_2, \dots, s_m\}$ over a finite alphabet
$\Sigma$, the problem is to find a common superstring for the set $S$ of the
minimum possible length. A string $s$ is considered a superstring of $S$ if
every string $s_i \in S$ appears as a contiguous substring (a consecutive block)
of $s$. The length of this optimal shortest superstring is denoted as $opt(S)$.
\end{problem}

Since the Shortest Common Superstring problem is NP-hard \cite{gallant1980finding}, no polynomial-time algorithm for finding an exact solution is known, and it is widely believed that none exists (assuming $\text{P} \neq \text{NP}$) \cite{garey1979computers}. Consequently, all known exact algorithms require exponential time in the worst case, making them impractical for large problem instances encountered in practice. Therefore, research has focused on the development of polynomial-time approximation algorithms that produce near-optimal solutions with provable approximation guarantees \cite{vazirani2001approximation}.\\
The SCS problem has numerous real-world applications across multiple scientific
and engineering fields:
\begin{itemize}
    \item \textbf{Bioinformatics:} One of the most prominent applications of the
    Shortest Common Superstring problem is DNA sequence assembly
    \cite{setubal1997introduction}. In shotgun sequencing, long DNA
    molecules—often consisting of millions of base pairs—are fragmented into
    thousands of short, overlapping reads, typically around 500 bases long. The
    original DNA sequence is then reconstructed by finding a short superstring
    that contains all reads, making SCS a natural computational model for this
    task \cite{peltola1984seqa}.
    \item \textbf{Data Compression:} The SCS serves as a theoretical foundation
    for dictionary-based data compression \cite{storer1982data}. A common
    superstring can be viewed as a maximally compressed form of a set of
    keywords, drastically decreasing data redundancy.
    \item \textbf{Operations Research:} Beyond strings and text, SCS solutions
    can be abstracted to solve scheduling problems, specifically in routing or
    scheduling operations on machines that require coordinated starting times
    \cite{lawler1985traveling}.
\end{itemize}
\begin{table}[htpb]
    \centering
    \caption{History of approximation-ratio improvements for the Shortest Common
    Superstring (SCS) problem. Rows in \textbf{bold} indicate the algorithm(s)
    discussed in this work.}
    \label{tab:scs_approximation_history}
    \renewcommand{\arraystretch}{1.3}
    \resizebox{\textwidth}{!}{%
    \begin{tabular}{@{}lllc@{}}
        \toprule
        \textbf{Year} & \textbf{Authors} & \textbf{Algorithm / Main Approach} & \textbf{Approximation Ratio} \\
        \midrule
        1994 & Blum et al. \cite{blum1994linear} & Cycle Cover / Modified Greedy & $3$ \\
        1997 & Teng, Yao \cite{teng1997approximating} & Optimal Cycle Covers / Matchings & $2 \frac{8}{9} \approx 2.89$ \\
        1997 & Czumaj et al. \cite{czumaj1997sequential} & Prefix Graph / Cycle Covers & $2 \frac{5}{6} \approx 2.83$ \\
        1998 & Armen, Stein \cite{armen19982} & Graph Matching Frameworks & $2 \frac{2}{3} \approx 2.67$ \\
        \textbf{1997} & \textbf{Breslauer, Jiang, Jiang \cite{Breslauer1996}} & \textbf{String Periodicity / Overlap Rotation (Algorithm 1)} & $\mathbf{2 \frac{2}{3} \approx 2.67}$ \\
        \textbf{1997} & \textbf{Breslauer, Jiang, Jiang \cite{Breslauer1996}} & \textbf{String Periodicity / Overlap Rotation (Algorithm 2)} & $\mathbf{2 \frac{25}{42} \approx 2.596}$ \\
        2005 & Kaplan et al. \cite{kaplan2005approximation} & ATSP Cycle Cover Reduction & $2 \frac{1}{2} = 2.50$ \\
        2013 & Mucha \cite{mucha2013lyndon} & ATSP Bounding via Lyndon Words & $2 \frac{11}{23} \approx 2.478$ \\
        2023 & Englert et al. \cite{englert2023approximation} & Improved Cycle Overlap Bounds & $\frac{\sqrt{67}+14}{9} \approx 2.466$ \\
        \bottomrule
    \end{tabular}%
    }
\end{table}

\section{Distance Graph and Cycle Covers}

It turns out that there is a relation between the Shortest Common Superstring
problem and certain graph problems, which forms the basis of the approximation
algorithms discussed in this work.

\begin{definition}[Traveling Salesperson Problem (TSP) \cite{karp1972reducibility}]
Given a complete directed weighted graph $G = (V, E, w)$, the Traveling
Salesperson Problem (TSP) asks for a minimum weight Hamiltonian cycle in $G$,
i.e., a minimum weight cycle visiting every vertex of $G$ exactly once.
\end{definition}

\begin{definition}[Distance Graph \cite{vazirani2001approximation}]
The distance graph for a set of strings $S = \{s_1, \dots, s_m\}$ is a directed
weighted graph $G_S = (V,E,w)$, in which the set of vertices is $V = S$, the set
of edges is $E = \{(s_i, s_j) \colon 1 \le i \neq j \le m\}$, and the weight
function assigns to each edge the distance between the corresponding strings,
i.e., $w(s_i, s_j) = d(s_i, s_j)$.
\end{definition}

If we denote the cost of a minimum weight Hamiltonian cycle on $G_S$ as
$TSP(G_S)$, it provides a very good approximation of the shortest superstring.
Specifically, for any $s_i \in S$, the relationship $TSP(G_S) \le opt(S) \le
TSP(G_S) + |s_i|$ holds. However, since the TSP is NP-hard
\cite{karp1972reducibility}, approximation algorithms utilize a relaxed version
of this problem known as the cycle cover problem (also referred to as the
assignment problem).

\begin{definition}[Cycle Cover \cite{vazirani2001approximation}]
A cycle cover in a directed weighted graph is a set of simple cycles such that
every vertex of the graph is contained in exactly one of them. The weight of the
cover is the total sum of the weights of all cycles comprising it.
\end{definition}

Unlike the TSP, a minimum weight cycle cover on any directed weighted graph can
be efficiently computed in polynomial time. Specifically, it can be solved in
$O(n^3)$ time using the Hungarian algorithm \cite{kuhn1955hungarian, edmondskarp1972, tomizawa1971techniques}.

Let $CYC(G_S)$ denote the weight of a minimum weight cycle cover of $G_S$. This
weight provides a useful lower bound for the length of the optimal superstring,
establishing the inequality:
\begin{equation}\label{eq:1.3}
CYC(G_S) \le TSP(G_S) \le opt(S)
\end{equation}
However, because there is no obvious general upper bound on $opt(S)$ expressed
solely in terms of $CYC(G_S)$, designing effective approximation algorithms
relies on specific properties of the strings. In fact, on general distance
graphs the gap between $CYC(G_S)$ and $TSP(G_S)$ can be made arbitrarily large
\cite{sahni1976pcomplete}.
